\par\Needspace{4\baselineskip}
\begin{center}
\hypertarget{native-056-article038}{}\pdfbookmark[2]{38.}{article-038}38.
\end{center}
Es ist hilfreich, das im vorigen Artikel Dargelegte an einem Beispiel zu erläutern. Für das Beispiel seien die wahre Anomalie $=100^0$ und die Exzentrizität $=0{,}967\,6456\,7$ gegeben. Außerdem gelte $\log q=9{,}765\,6500$. Die Rechnung für $E$, $B$, $A$ und $t$ sieht dann so aus.
\begin{center}
\begin{tabular}{@{}ll@{}}
\makebox[18em][l]{$\log\operatorname{tang}\tfrac{1}{2}v$\hfill .\,.\,.\,.} & $0{,}076\,1865$\\[.3em]
\makebox[18em][l]{$\log\BareRoot{\dfrac{1-e}{1+e}}$\hfill .\,.\,.\,.\,.} & $9{,}107\,9927$\\[.4em]
\hline
\makebox[18em][l]{$\log\operatorname{tang}\tfrac{1}{2}E$\hfill .\,.\,.} & $9{,}184\,1792$,
\end{tabular}
\end{center}
\noindent Daraus folgen $\tfrac{1}{2}E=8^0\,41^{\prime}\,19^{\prime\prime}{,}32$ und $E=17^0\,22^{\prime}\,38^{\prime\prime}{,}64$. Diesem Wert von $E$ entspricht $\log B=0{,}000\,0040$. Ferner findet man in Teilen des Radius $E=0{,}303\,2928$ und $\sin E=0{,}298\,6643$. Daraus ergibt sich $\tfrac{9}{20}E+\tfrac{1}{20}\sin E=0{,}151\,4150$ mit dem Logarithmus $=9{,}180\,1689$ und folglich $\log A^{\frac{1}{2}}=9{,}180\,1649$. Nach Formel 1) des vorigen Artikels erhält man die folgende Rechnung.
\begin{center}
\begin{tabular}{@{}lrl l@{}}
\multicolumn{3}{@{}l}{\makebox[22em][l]{$\log\dfrac{2Bq^{\frac{3}{2}}}{k\BareRoot{(1-e)}}$\hfill .\,.\,.\,.\,.\,.\,.\,.\,.}} & $2{,}458\,9614$\\[1.0em]
\multicolumn{3}{@{}l}{\makebox[22em][l]{$\log A^{\frac{1}{2}}$\hfill .\,.\,.\,.\,.\,.\,.\,.\,.\,.\,.\,.}} & $9{,}180\,1649$\\[.3em]
\hline
$\log$ & $43{,}56386$ & $=$ & $1{,}639\,1263$\\[.3em]
 & $19{,}98014$ & & \\
\cline{2-2}
 & $63{,}54400$ & $=t$ &
\end{tabular}
\end{center}
\begin{center}
\begin{tabular}{@{}ll@{}}
\makebox[22em][l]{$\log\dfrac{2B(1+9e)}{15k}\left(\dfrac{q}{1-e}\right)^{\frac{3}{2}}$\hfill .\,.\,.} & $3{,}760\,1038$\\[.4em]
\makebox[22em][l]{$\log A^{\frac{3}{2}}$\hfill .\,.\,.\,.\,.\,.\,.\,.\,.\,.\,.\,.} & $7{,}540\,4947$\\[.3em]
\hline
$\log 19{,}98014$\hfill $=$ & $1{,}300\,5985$
\end{tabular}
\end{center}
\noindent Behandelt man dasselbe Beispiel nach dem gewöhnlichen Verfahren, so ergibt sich für $e\sin E$ in \SourceLeaf{057}Bogensekunden $=59610^{\prime\prime}{,}79=16^0\,33^{\prime}\,30^{\prime\prime}{,}79$ und damit für die mittlere Anomalie $=49^{\prime}\,7^{\prime\prime}{,}85=2947^{\prime\prime}{,}85$. Daraus und aus $\log k\left(\dfrac{1-e}{q}\right)^{\frac{3}{2}}=1{,}666\,4302$ folgt $t=63{,}54410$. Die Differenz, die hier nur den $10\,000$. Teil eines Tages beträgt, hätte leicht drei- oder viermal so groß ausfallen können, wenn die Fehler in dieselbe Richtung gewirkt hätten.

Im Übrigen ist klar, dass eine solche Tafel für $\log B$ allein genügt, um auch das inverse Problem mit jeder gewünschten Genauigkeit zu lösen. Man bestimmt $E$ durch wiederholtes Probieren so, dass der daraus berechnete Wert von $t$ mit dem vorgegebenen übereinstimmt. Dieses Verfahren wäre jedoch recht mühsam. Wir zeigen daher jetzt, wie man die Hilfstafel viel zweckmäßiger einrichten, planloses Probieren ganz vermeiden und die gesamte Rechnung auf einen übersichtlichen und raschen Algorithmus zurückführen kann, der nichts zu wünschen übrig zu lassen scheint.
